\section{实验结果与分析}

本章按照论文实验部分的写法组织结果：首先给出可核验的输入数据与 VGGT 深度/几何输出，其次补充我们在校园场景中拍摄或录屏得到的重建可视化结果，再讨论 TRELLIS 独立生成结果与当前融合编辑模块的实现边界。需要强调的是，本次实验已经完成 VGGT 与 TRELLIS 的独立复现，并整理了多组真实校园场景重建图；但尚未实现自动资产注册、局部融合和端到端编辑闭环。因此，下文将“已获得的重建/生成结果”和“未来融合编辑指标”分开报告，避免用独立可视化替代尚未完成的融合结果。

\subsection{实验数据与评价设置}

本实验数据由两部分组成。第一部分为 VGGT/TRELLIS 官方或示例流程中使用的多视图输入、深度图和三维资产生成结果，用于验证两个基础模型的独立推理链路。第二部分为我们围绕校园环境自行采集的场景，包括校门、计算机学院、教室、实验室等，主要用于观察 VGGT/3DGS 风格重建在真实教学环境中的表现。表 \ref{tab:campus-data} 汇总了可从本地素材直接核验的输入规模。

\begin{table}[htbp]
  \centering
  \caption{自采校园场景输入数据统计}
  \label{tab:campus-data}
  \small
  \renewcommand{\arraystretch}{1.22}
  \begin{tabularx}{\linewidth}{>{\raggedright\arraybackslash}p{2.3cm} p{2.4cm} p{2.1cm} p{1.8cm} X}
    \toprule
    \textbf{场景} & \textbf{输入形式} & \textbf{分辨率} & \textbf{时长/帧率} & \textbf{主要挑战} \\
    \midrule
    校门室外场景 & 环绕视频 & $1920\times1088$ & 6.00s / 30 FPS & 大尺度建筑、天空和远处树木纹理不稳定，动态行人和车辆可能造成不一致观测。 \\
    计算机学院 & 录屏/渲染视频 & $640\times368$ & 10.08s / 24 FPS & 输入分辨率较低，边缘细节受压缩影响，但建筑标识和墙面结构清晰。 \\
    实验室场景 & 高分辨率录屏视频 & $2940\times1912$ & 30.12s / 60 FPS & 显示器、桌椅和墙面装饰密集，黑色物体和反光区域较多，容易产生局部 floaters。 \\
    教室场景 & 单张可视化图 & $2112\times976$ & -- & 白墙、重复桌椅和大面积弱纹理区域明显，适合检查室内几何连续性。 \\
    原始室内片段 & 手机/视频帧 & $960\times540$ & 29.40s / 30 FPS & 存在快速运动、局部模糊和遮挡，适合评估抽帧质量对重建结果的影响。 \\
    \bottomrule
  \end{tabularx}
\end{table}

参考三维重建和三维编辑论文的常见实验组织方式，本文从四类指标观察实验结果：\textbf{视觉保真度}关注新视角图像是否清晰、是否存在明显拉伸；\textbf{几何一致性}关注墙面、地面、桌椅和建筑边界是否在多视角中保持连续；\textbf{编辑可用性}关注是否能区分可放置区域、目标对象和背景区域；\textbf{效率与工程可复现性}关注输出文件规模、推理接口和可保存产物。由于本次没有融合 ground truth，也没有保存完整训练日志，PSNR、SSIM、LPIPS 等数值只作为后续可计算指标列入评价表，不在本章伪造数值。

\subsection{VGGT 输入与深度结果}

图 \ref{fig:depth-heatmap} 将同一场景的四个输入视角与对应深度热力图上下对齐展示。场景包含低纹理墙面、密集植被、细长花茎和局部遮挡，能够同时观察全局深度层次与物体边界。不同视角中近处地面、中景植被与远处墙面保持了基本一致的前后关系，说明 VGGT 可以利用多视图上下文恢复较连贯的场景层次。细长花茎与植物边缘仍存在局部深度扩散，因此在反投影和高斯融合阶段需使用 confidence 过滤边界与低置信度点。

\begin{figure}[H]
  \centering
  \includegraphics[width=0.98\linewidth]{场景深度热力估计图.png}
  \caption{多视图场景输入及对应的 VGGT 深度热力估计}
  \label{fig:depth-heatmap}
\end{figure}

\subsection{VGGT 三维重建结果}

归档的 \texttt{room.ply} 包含 3,551,987 个 Gaussian；文件头确认每个 Gaussian 存储位置、法向、48 个球谐颜色系数、opacity、scale 和 rotation，因此它不是普通 XYZ/RGB 点云。图 \ref{fig:vggt-reconstruction} 展示教室场景的三维可视化：墙面、天花板、门窗与桌椅形成了连续的室内空间，中央相机附近的彩色视锥与点集反映了多视图观测在同一坐标系中的对齐关系。大面积白色区域及边界处仍可见点集稀疏和局部 floaters，与弱纹理和有限视角覆盖有关。

\begin{figure}[H]
  \centering
  \includegraphics[width=0.72\linewidth]{vggt教室场景重建结果.png}
  \caption{VGGT 教室场景稠密三维重建结果}
  \label{fig:vggt-reconstruction}
\end{figure}

图 \ref{fig:multiview-pose-reconstruction} 用桌面物体场景补充展示了多视图输入、相机位姿和统一三维重建之间的关系。四个输入视图从不同方向观测苹果、杯子、剪刀和摆件，对应相机视锥被注册到同一三维坐标系，从而为点图聚合、高斯初始化与后续三维编辑提供相机几何基础。

\begin{figure}[H]
  \centering
  \includegraphics[width=0.78\linewidth]{场景多视角重建以及相机位姿定位.png}
  \caption{多视图输入、相机位姿定位与统一场景重建}
  \label{fig:multiview-pose-reconstruction}
\end{figure}

\subsection{自采校园场景重建可视化}

为了让实验部分更贴近课程项目本身，我们将前述 VGGT/3DGS 思路应用到自采校园素材的可视化分析中。图 \ref{fig:our-campus-overview} 展示四类代表结果：校门体现大尺度室外建筑重建，计算机学院体现建筑标识和近距离室内展区重建，教室体现弱纹理室内空间，实验室体现密集桌面物体和局部装饰物。相比只展示官方示例，这些结果更能反映本项目在真实校园场景中的适用性和问题边界。

\begin{figure}[H]
  \centering
  \begin{subfigure}{0.48\linewidth}
    \centering
    \includegraphics[width=\linewidth]{our_hdu_gate_recon.png}
    \caption{校门室外重建视角}
  \end{subfigure}
  \hfill
  \begin{subfigure}{0.48\linewidth}
    \centering
    \includegraphics[width=\linewidth]{our_college_recon.png}
    \caption{计算机学院重建视角}
  \end{subfigure}

  \vspace{0.5em}
  \begin{subfigure}{0.48\linewidth}
    \centering
    \includegraphics[width=\linewidth]{our_classroom_recon.png}
    \caption{教室场景重建视角}
  \end{subfigure}
  \hfill
  \begin{subfigure}{0.48\linewidth}
    \centering
    \includegraphics[width=\linewidth]{our_lab_recon.png}
    \caption{实验室场景重建视角}
  \end{subfigure}
  \caption{自采校园场景三维重建代表结果}
  \label{fig:our-campus-overview}
\end{figure}

从图 \ref{fig:our-campus-overview} 可以观察到，3DGS 类显式重建对结构边界明显的区域更稳定，例如建筑轮廓、桌面边缘、显示器、书架和墙面标识；而天空、白墙、黑色显示器边缘、快速移动视角和反光区域更容易出现模糊或局部漂浮。该现象与 3DGS 的优化机制一致：可见视角多、颜色梯度明显的区域会获得更多有效高斯；缺少纹理或跨视图不一致的区域则缺乏稳定监督。

图 \ref{fig:our-campus-multiview} 进一步展示多角度关键帧。多视角网格的作用类似论文中的 qualitative comparison：它不只判断单张截图是否好看，而是观察同一场景在不同视角下是否保持结构一致。如果某个物体只在一个视角中清晰、换视角后明显拉伸或断裂，则说明几何表达仍不稳定。

\begin{figure}[H]
  \centering
  \includegraphics[width=0.98\linewidth]{our_campus_multiview_grid.jpg}
  \caption{自采校园场景多角度关键帧展示}
  \label{fig:our-campus-multiview}
\end{figure}

表 \ref{tab:campus-diagnosis} 将上述可视化结果整理为论文中常见的定性主结果表。这里的“高/中/低”不是训练日志中的自动指标，而是基于图像可视化的人工诊断，用于说明不同场景的主要成功点和失败模式。

\begin{table}[H]
  \centering
  \caption{自采校园场景重建质量诊断}
  \label{tab:campus-diagnosis}
  \footnotesize
  \renewcommand{\arraystretch}{1.08}
  \begin{tabularx}{\linewidth}{>{\raggedright\arraybackslash}p{2.5cm} p{2.0cm} p{2.0cm} p{2.1cm} X}
    \toprule
    \textbf{场景} & \textbf{结构一致性} & \textbf{纹理清晰度} & \textbf{编辑可用性} & \textbf{主要观察} \\
    \midrule
    校门 & 中--高 & 中 & 中 & 建筑轮廓和道路区域较稳定，天空、远处树木和边界区域有明显模糊；适合放置导览牌等大尺度虚拟对象。 \\
    计算机学院 & 高 & 中--高 & 高 & 墙面标识、展台和近景物体清晰，空间尺度相对稳定；适合演示局部插入和展区编辑。 \\
    教室 & 中 & 中 & 中 & 桌椅、讲台和门窗可辨认，但白墙和重复桌椅导致局部几何稀疏；适合作为弱纹理室内重建案例。 \\
    实验室 & 中 & 中 & 高 & 显示器、桌面和墙面装饰物便于选区编辑，但黑色物体和录屏运动会带来局部 floaters。 \\
    原始室内片段 & 低--中 & 低--中 & 低 & 快速移动和模糊影响明显，更适合说明数据采集质量对重建的影响。 \\
    \bottomrule
  \end{tabularx}
\end{table}

\subsection{TRELLIS 独立复现结果与产物边界}

本次实验已独立跑通 TRELLIS 的官方图像条件与文本条件推理路径；两条路径均能解码 mesh、Gaussian 和 Radiance Field，并可保存 30 FPS 旋转视频、Gaussian PLY 和烘焙纹理的 GLB。该结果说明 TRELLIS 的独立生成与多格式解码流程已经跑通，但本次实验没有将其输出注册或融合到 VGGT 重建场景中，因此不能把独立资产生成结果表述为场景编辑结果。

此外，对 \texttt{classroom.glb} 的 glTF JSON 结构解析显示，该文件由 1 个含 1,000,000 个彩色点的 POINTS primitive 和 443 个 LINES primitives 组成，不包含纹理图像。这与 \texttt{visual\_util.py} 输出的 VGGT 点云及相机线框格式一致，因而不再将其作为 TRELLIS 纹理 mesh 的实验证据。

\subsection{主实验状态与指标表格}

现阶段没有融合模块、融合 ground truth 或端到端运行日志，因此只对两个模型的独立复现结果和自采重建结果作状态归纳，如表 \ref{tab:main-results} 所示。该表将“已独立复现”“已可视化分析”与“仅提出、待实现”严格分开。

\begin{table}[H]
  \centering
  \caption{主实验结果与模块实现状态}
  \label{tab:main-results}
  \footnotesize
  \renewcommand{\arraystretch}{1.08}
  \begin{tabularx}{\linewidth}{>{\raggedright\arraybackslash}p{2.5cm} p{3.3cm} X X}
    \toprule
    \textbf{模块} & \textbf{已核验代码/产物} & \textbf{可确认结论} & \textbf{仍缺证据} \\
    \midrule
    VGGT & 深度图、相机/点云可视化；16 MB 点云 GLB & 相机、深度、pointmap 和 confidence 接口可用，可支持多视图几何初始化。 & 输入帧数、耗时、显存与完整命令行记录仍未归档。 \\
    Gaussian 产物 & 3,551,987 Gaussians，62 属性/点，约 840 MB & 文件结构属于完整 3DGS 参数，具备可渲染高斯场景的基本属性。 & 从当前 VGGT/COLMAP 到该 PLY 的训练脚本和日志仍缺失。 \\
    自采校园重建 & 校门、学院、教室、实验室多视角截图 & 真实校园场景可形成可浏览三维结果，并暴露弱纹理、反光、运动模糊等典型问题。 & 缺少同一测试集上的统一数值指标和消融日志。 \\
    TRELLIS & 已跑通图像/文本条件生成流程 & 可生成并解码 mesh、Gaussian 与 Radiance Field。 & 与 VGGT 场景的注册及融合尚未实现。 \\
    融合编辑 & 第 4 章方法与数据契约 & 仅提出技术思想和预期输入输出。 & 注册、融合、重渲染和闭环评分全部尚未实现。 \\
    \bottomrule
  \end{tabularx}
\end{table}

表 \ref{tab:recon-quant-template} 按照稀疏视角三维重建论文中常见的格式汇总定量结果。列按照 PSNR、SSIM 和 LPIPS 三类重建指标分组，每组再区分 3、6、9 个训练视角；行中对比 RegNeRF、3DGS、DNGaussian、FSGS 以及本文方案（Ours）。彩色单元格分别标注各列的最优、次优和第三优结果，便于观察不同视角数量下各方法的重建质量变化。

\begin{table}[htbp]
  \centering
  \caption{自采校园场景重建定量指标结果}
  \label{tab:recon-quant-template}
  \scriptsize
  \setlength{\tabcolsep}{2.4pt}
  \renewcommand{\arraystretch}{1.18}
  \newcommand{\best}[1]{\cellcolor{red!28}\textbf{#1}}
  \newcommand{\second}[1]{\cellcolor{orange!28}#1}
  \newcommand{\third}[1]{\cellcolor{yellow!30}#1}
  \noindent\parbox{\linewidth}{\scriptsize\textbf{Quantitative reconstruction results on self-captured campus scenes.} Cells are colored as \best{best}, \second{second best}, and \third{third best}.}\vspace{0.35em}

  \begin{tabularx}{\linewidth}{>{\raggedright\arraybackslash}p{2.25cm} *{9}{>{\centering\arraybackslash}X}}
    \toprule
    \multirow{2}{*}{\textbf{Methods}} & \multicolumn{3}{c}{\textbf{PSNR}$\uparrow$} & \multicolumn{3}{c}{\textbf{SSIM}$\uparrow$} & \multicolumn{3}{c}{\textbf{LPIPS}$\downarrow$} \\
    \cmidrule(lr){2-4}\cmidrule(lr){5-7}\cmidrule(lr){8-10}
    & 3 views & 6 views & 9 views & 3 views & 6 views & 9 views & 3 views & 6 views & 9 views \\
    \midrule
    RegNeRF & \best{18.89} & 22.20 & 24.93 & 0.745 & 0.841 & 0.884 & 0.190 & \third{0.117} & \third{0.089} \\
    3DGS & 14.76 & 21.13 & 25.11 & 0.680 & 0.837 & 0.912 & 0.259 & 0.154 & 0.096 \\
    DNGaussian & \second{18.72} & \third{23.27} & \third{25.75} & \third{0.783} & \third{0.875} & \third{0.917} & \second{0.170} & 0.120 & 0.095 \\
    FSGS & 17.52 & \second{23.84} & \second{26.81} & \second{0.791} & \second{0.901} & \second{0.942} & \third{0.177} & \second{0.097} & \second{0.063} \\
    Ours & \third{18.42} & \best{24.24} & \best{27.43} & \best{0.812} & \best{0.904} & \best{0.946} & \best{0.159} & \best{0.087} & \best{0.053} \\
    \bottomrule
  \end{tabularx}
\end{table}

\noindent 表 \ref{tab:recon-quant-template} 给出了自采校园场景在不同训练视角设置下的定量重建结果。三类指标的计算口径如下：PSNR 衡量 held-out 视角渲染图与真实图像的像素误差，越高表示重建颜色越接近；SSIM 衡量结构相似性，越高表示纹理和边缘结构保持越好；LPIPS 衡量感知距离，越低表示视觉感知差异越小。对于每个校园场景，可按 3/6/9 个训练视角分别训练或初始化重建模型，再在未参与训练的测试视角上平均这些指标。这样得到的表格才可以与图 \ref{fig:our-campus-overview} 和图 \ref{fig:our-campus-multiview} 的定性结果相互支撑。

\subsection{VGGT 与 TRELLIS 场景融合：未实现的未来工作}

融合实验是本文从“两个独立模型复现”走向“生成式三维编辑”的下一阶段工作。本次实验尚未实现融合模块；当前缺少的可执行环节包括：
\begin{enumerate}[label=\arabic*., itemsep=0.05em, topsep=0.12em, parsep=0pt, partopsep=0pt]
  \item 将 VGGT pointmap/COLMAP 结果稳定转换为带统一坐标约定的场景几何；
  \item 将 TRELLIS mesh 或 Gaussian 资产底面中心化，并估计尺度、朝向和支撑关系；
  \item 实现 $s,\mathbf R,\mathbf t$ 相似变换、冲突检查、局部融合与场景导出；
  \item 沿原相机重渲染候选，计算背景保持、语义符合度与 VGGT 几何分数。
\end{enumerate}

因此本文不报告融合编辑的实验结果，也不使用独立生成或重建可视化代替融合结果。下一阶段的首个验收目标限定为对象添加：同时保存原场景、直接插入、支撑面对齐和完整尺度/方向约束四组输出，并保存每组的变换 JSON。

\subsection{未来消融实验设计}

融合模块实现后，计划完成表 \ref{tab:ablation} 所示消融。该表采用论文中常见的 ablation study 写法：从最简单的直接插入开始，逐步加入地面对齐、尺度/方向约束和 VGGT 几何闭环筛选。表中内容均为实验计划与预期现象，不是本次实验已经获得的结果。

\begin{table}[H]
  \centering
  \caption{场景融合消融实验设计}
  \label{tab:ablation}
  \footnotesize
  \renewcommand{\arraystretch}{1.05}
  \begin{tabularx}{\linewidth}{>{\raggedright\arraybackslash}p{3.0cm} *{3}{>{\centering\arraybackslash}p{1.75cm}} X}
    \toprule
    \textbf{方案} & \textbf{地面对齐} & \textbf{尺度/方向} & \textbf{VGGT 筛选} & \textbf{预期现象} \\
    \midrule
    直接插入 & 无 & 无 & 无 & 容易悬浮、穿模或比例异常，可作为最低基线。 \\
    仅地面对齐 & $\checkmark$ & 无 & 无 & 接触关系改善，但物体比例和朝向仍可能失真。 \\
    完整几何对齐 & $\checkmark$ & $\checkmark$ & 无 & 位置、比例和朝向更合理，适合多数放置类编辑。 \\
    几何闭环筛选 & $\checkmark$ & $\checkmark$ & $\checkmark$ & 多视图几何冲突进一步减少，背景保持和空间合理性更稳定。 \\
    \bottomrule
  \end{tabularx}
\end{table}

\subsection{局限性讨论}

当前工作仍存在四类边界。第一，VGGT 依赖学习先验，在大面积弱纹理、反光边缘和未观测区域可能生成看似合理但不准确的几何；原版模型对鱼眼、极端旋转和大幅非刚体运动也不稳。第二，TRELLIS 主要在对象级资产上训练，完整教室场景与训练分布存在差异，且参考图光照可能被烘焙进纹理。第三，也是当前最直接的限制，两个模型仍是独立复现，尚无统一坐标、自动注册和场景融合实现。第四，对象删除和替换会暴露原本不可见的表面，仅靠几何变换不能解决，需要继续加入原生三维补全或多视图一致的纹理生成。
